\documentclass[12pt]{article}

%% Basic document formatting
\usepackage{amsmath, amsthm, amssymb, amsfonts}
\usepackage{mathtools}
\usepackage{xspace}
\usepackage{thmtools}
\usepackage{graphicx}
\usepackage{tikz}
\usepackage{ytableau}
\usepackage{blkarray}
\usepackage{hyperref}
\usepackage{setspace}
\usepackage{array}
\usepackage{fancyhdr}
\usepackage{titling}
\usepackage[left=0.4in,right=0.4in,top=1in,bottom=1in]{geometry}
\usepackage{float}
\usepackage{cancel}
\usepackage{tabularx}
\usepackage[utf8]{inputenc}
\usepackage[english]{babel}
\usepackage{framed}
\usepackage[dvipsnames]{xcolor}
\usepackage{environ}
\usepackage{tcolorbox}
\tcbuselibrary{theorems,skins,breakable}

\usepackage{enumitem}
\setlist[enumerate]{leftmargin=*}
\setlist[enumerate,1]{labelindent=\parindent}
\setlist[enumerate,2]{labelindent=0pt}

% Blackboard Bold
\newcommand{\Z}{\mathbb{Z}}                 % integers
\newcommand{\Zpl}{\mathbb{Z}_{+}}           % positive integers
\newcommand{\Znn}{\mathbb{Z}_{\geq 0}}      % nonnegative integers. 
\newcommand{\Q}{\mathbb{Q}}                 % rationals
\newcommand{\Qpl}{\mathbb{Q}_{+}}           % positive Rationals
\newcommand{\R}{\mathbb{R}}                 % reals
\newcommand{\Rpl}{\mathbb{R}_{+}}           % positive Reals
\newcommand{\C}{\mathbb{C}}                 % complex numbers
\newcommand{\F}{\mathbb{F}}                 % field

% Words
\newcommand{\st}{\text{ such that }}
\newcommand{\wrt}{\text{ with respect to }}
\newcommand{\with}{\text{ with }}
\newcommand{\ie}{\text{, i.e. }}
\newcommand*{\ditto}{\text{---}\texttt{"}\text{---}}

% Operators
\newcommand{\abs}[1]{\left|#1\right|}                   % absolute value
\newcommand{\norm}[1]{\abs{\abs{#1}}}
\newcommand{\floor}[1]{\left\lfloor #1 \right\rfloor}   % floor
\newcommand{\ceil}[1]{\left\lceil #1 \right\rceil}      % ceiling
\newcommand{\powset}[1]{\mathscr{P}(#1)}

% Category Theory 
\renewcommand{\hom}{\text{Hom}}
\newcommand{\homspace}[3]{\hom_{#1}(#2, #3)}
\newcommand{\coproduct}[9]{
  \begin{tikzcd}[ampersand replacement=\&]
      \& \& #1 \arrow[dll, bend right, "#6"] \arrow[dl, "#5"] \\
   #4 \& #3 \arrow[l, "#9"] \\
      \& \& #2 \arrow[ull, bend left, "#8"] \arrow[ul, "#7"]
  \end{tikzcd}
}

\newcommand{\product}[9]{ 
  \begin{tikzcd}[ampersand replacement=\&]
      \& \& #3 \\
      #1 \arrow[r, "#7"] \arrow[urr, bend left, "#5"] \arrow[drr, bend right, "#6"] \& #2 \arrow[ur, "#8"] \arrow[dr, "#9"] \\ 
      \& \& #4
  \end{tikzcd}
}


% Algebra 
\newcommand{\Syl}[2]{\text{Syl}_{#1}{(#2)}}           % set of Sylow-p subgroups 
\newcommand{\<}{\langle}                % \<x\>, subgroup generated by x 
\renewcommand{\>}{\rangle}
\renewcommand{\index}[2]{[#1 : #2]}
\newcommand{\id}{\text{id}}             % identity element
\newcommand{\lhdneq}{%
  \mathrel{\ooalign{$\lneq$\cr\raise.22ex\hbox{$\lhd$}\cr}}}
\newcommand{\order}[1]{\text{o}(#1)}    % order of an element
\newcommand{\spec}{\operatorname{Spec}}
\newcommand{\rad}{\operatorname{rad}}   % radical of an ideal 
\newcommand{\sym}[1]{\mathfrak{S}_{#1}}
\newcommand{\aut}[1]{\text{Aut}(#1)} 
\newcommand{\gal}[2]{\text{Gal}(#1 / #2)} 
\newcommand{\inn}[1]{\text{Inn}(#1)} 
\let\oldcong\cong
\let\oldequiv\equiv
\renewcommand{\cong}{\oldequiv}
\renewcommand{\equiv}{\oldcong}

\newcommand{\triangleleftneq}{\mathrel{\ooalign{%
  $\trianglelefteq$\cr
  \hidewidth\raise0.06ex\hbox{$\not\mkern4mu$}\hidewidth\cr
}}}

\makeatletter % cycle
\newcommand{\cyc}[1]{(\mathbf{\cyc@process#1\relax})}
\def\cyc@process#1#2\relax{%
	#1%
	\ifx\relax#2\relax
	\else
		\,\cyc@process#2\relax
	\fi
}
\makeatother

\newcommand{\GL}[2]{\text{GL}_{#1}(#2)}                             % general linear group
\newcommand{\SL}[2]{\text{SL}_{#1}(#2)}                             % special linear group
\newcommand{\gl}[2]{\mathfrak{gl}_{#1}(#2)}
\renewcommand{\sl}[2]{\mathfrak{sl}_{#1}(#2)}
\newcommand{\so}[2]{\mathfrak{so}_{#1}(#2)}
\newcommand{\su}[1]{\mathfrak{su}(#1)}

% Linear Algebra
\newcommand{\bmat}[1]{\begin{bmatrix}#1\end{bmatrix}}   % bracketed matrix
\newcommand{\rank}{\operatorname{rank}}                 % rank
\newcommand{\nullity}{\operatorname{nullity}}           % nullity
\newcommand{\tr}{\operatorname{tr}}

% Combinatorics
\newcommand{\qnum}[1]{\left[#1\right]_q}                % q-analog
\newcommand{\qfact}[1]{\left[#1\right]!_q}              % q-analog factorial 
\newcommand{\qbinom}[2]{\binom{#1}{#2}_q}               % Gaussian binomial coefficient

% Topology/Analysis
\newcommand{\ball}[2]{\text{B}_{#1}(#2)}  % B_r(x): open r-balls around x
\newcommand{\diam}{\text{diam}}           % diamter of a set in metric space 
\newcommand{\lowint}[2]{\underline{\int_{#1}^{#2}}}
\newcommand{\upint}[2]{\overline{\int}_{#1}^{#2}}

% Misc Notation
\newcommand{\oneton}[1]{\{1, \ldots, #1\}}
\newcommand{\defeq}{\vcentcolon=}   % :=
\newcommand{\eqdef}{=\vcentcolon}   % =: 
\renewcommand{\bf}[1]{\textbf{#1}}
\newcommand{\im}{\operatorname{im}}
\newcommand{\fnrest}[2]{\left.#1\right|_{#2}}
\newcommand{\isom}{\overset{\sim}{\rightarrow}} 


\renewcommand{\thesection}{\arabic{section}.}
\renewcommand{\thesubsection}{(\alph{subsection})}

\newtheorem{claim}{Claim}
\newtheorem*{lemma}{Lemma}

%% Headers & title setup
\newcommand{\course}{Math188}
\newcommand{\myname}{Jay Ser}
\setlength{\headheight}{14.5pt}
\pagestyle{fancy}
\fancyhf{}
\renewcommand{\headrulewidth}{0.4pt}
\lhead{\course}
\rhead{\myname}
\cfoot{\thepage}
\setlength{\droptitle}{-4em} 
\title{\course\ - HW \#6}
\author{\myname}
\date{2026.03.09}

\begin{document}
\maketitle
\thispagestyle{fancy}

%------------------------------------------------------------------------------%

\section{} 
It follows immediately from the property of the M\"{o}bius function that 
$$\mu_{P \times Q}[(x, y), (x, y)] = 1 = \mu_{P}[x, x] \mu_{Q}[y, y]$$
Suppose $(x, y) < (x', y')$.  
Then 
\begin{align*}
  \sum_{(x, y) \leq (z, w) \leq (x', y')} \mu_{P}[x, z] \mu_{Q}[y, w] 
  & = \sum_{x \leq z \leq x'} \mu_{P}[x, z] \sum_{y \leq w \leq y'} \mu_{Q}[y, w]  \\ 
  & = 0 
\end{align*}
since $x < x'$ or $y < y'$.
$\mu_{P \times Q}[(x, y), (x', y')]$ and $\mu_{P}[x, x'] \mu_{Q}[y, y']$ satisfy the same recursion, so they are equal. 

\section{} 
\begin{center}
\begin{tikzpicture}[
    every node/.style={
        circle, draw, inner sep=1pt, minimum size=24pt,
        font=\footnotesize\strut
    },
    edge/.style={thick, gray!60},
    rank0/.style={fill=orange!15, draw=orange!70!black},
    rank1/.style={fill=violet!12, draw=violet!70!black},
    rank2/.style={fill=pink!18, draw=pink!70!black},
]

% --- Rank 0 (bottom): discrete partition ---
\node[rank0] (0) at (0, 0) {$1|2|3$};

% --- Rank 1: three partitions with one pair merged ---
\node[rank1] (12) at (-2.5, 3) {$12|3$};
\node[rank1] (13) at (0, 3)   {$13|2$};
\node[rank1] (23) at (2.5, 3) {$23|1$};

% --- Rank 2 (top): single block ---
\node[rank2] (top) at (0, 6) {$123$};

% --- Rank labels ---
\node[draw=none, fill=none, font=\scriptsize\itshape\color{gray}] at (-4.5, 0) {rank 0};
\node[draw=none, fill=none, font=\scriptsize\itshape\color{gray}] at (-4.5, 3) {rank 1};
\node[draw=none, fill=none, font=\scriptsize\itshape\color{gray}] at (-4.5, 6) {rank 2};

% ========== EDGES ==========

% Rank 0 -> Rank 1
\draw[edge] (0) -- (12);
\draw[edge] (0) -- (13);
\draw[edge] (0) -- (23);

% Rank 1 -> Rank 2
\draw[edge] (12) -- (top);
\draw[edge] (13) -- (top);
\draw[edge] (23) -- (top);
\end{tikzpicture}
\end{center}

\begin{center}
\begin{tikzpicture}[
    node distance=1.5cm,
    every node/.style={
        circle, draw, inner sep=1pt, minimum size=22pt,
        font=\footnotesize\strut
    },
    edge/.style={thick, gray!60},
    rank0/.style={fill=orange!15, draw=orange!70!black},
    rank1/.style={fill=violet!12, draw=violet!70!black},
    rank2/.style={fill=green!12, draw=green!50!black},
    rank3/.style={fill=pink!18, draw=pink!70!black},
]

% --- Rank 0 (bottom): discrete partition ---
\node[rank0] (0) at (0, 0) {$1|2|3|4$};

% --- Rank 1: six partitions with one pair merged ---
\node[rank1] (12) at (-5, 3) {$12|3|4$};
\node[rank1] (13) at (-3, 3) {$13|2|4$};
\node[rank1] (14) at (-1, 3) {$14|2|3$};
\node[rank1] (23) at (1, 3)  {$23|1|4$};
\node[rank1] (24) at (3, 3)  {$24|1|3$};
\node[rank1] (34) at (5, 3)  {$34|1|2$};

% --- Rank 2: seven partitions with two blocks ---
\node[rank2] (123) at (-5.5, 6) {$123|4$};
\node[rank2] (124) at (-3.5, 6) {$124|3$};
\node[rank2] (134) at (-1.5, 6) {$134|2$};
\node[rank2] (234) at (0.5, 6)  {$234|1$};
\node[rank2] (1234a) at (2.5, 6) {$12|34$};
\node[rank2] (1324) at (4.2, 6)  {$13|24$};
\node[rank2] (1423) at (5.9, 6) {$14|23$};

% --- Rank 3 (top): single block ---
\node[rank3] (top) at (0, 9) {$1234$};

% --- Rank labels ---
\node[draw=none, fill=none, font=\scriptsize\itshape\color{gray}] at (-7.2, 0) {rank 0};
\node[draw=none, fill=none, font=\scriptsize\itshape\color{gray}] at (-7.2, 3) {rank 1};
\node[draw=none, fill=none, font=\scriptsize\itshape\color{gray}] at (-7.2, 6) {rank 2};
\node[draw=none, fill=none, font=\scriptsize\itshape\color{gray}] at (-7.2, 9) {rank 3};

% ========== EDGES ==========

% Rank 0 -> Rank 1
\foreach \x in {12,13,14,23,24,34} {
    \draw[edge] (0) -- (\x);
}

% Rank 1 -> Rank 2
% 12|3|4 covers: 123|4, 124|3, 12|34
\draw[edge] (12) -- (123);
\draw[edge] (12) -- (124);
\draw[edge] (12) -- (1234a);

% 13|2|4 covers: 123|4, 134|2, 13|24
\draw[edge] (13) -- (123);
\draw[edge] (13) -- (134);
\draw[edge] (13) -- (1324);

% 14|2|3 covers: 124|3, 134|2, 14|23
\draw[edge] (14) -- (124);
\draw[edge] (14) -- (134);
\draw[edge] (14) -- (1423);

% 23|1|4 covers: 123|4, 234|1, 14|23
\draw[edge] (23) -- (123);
\draw[edge] (23) -- (234);
\draw[edge] (23) -- (1423);

% 24|1|3 covers: 124|3, 234|1, 13|24
\draw[edge] (24) -- (124);
\draw[edge] (24) -- (234);
\draw[edge] (24) -- (1324);

% 34|1|2 covers: 134|2, 234|1, 12|34
\draw[edge] (34) -- (134);
\draw[edge] (34) -- (234);
\draw[edge] (34) -- (1234a);

% Rank 2 -> Rank 3
\foreach \x in {123,124,134,234,1234a,1324,1423} {
    \draw[edge] (\x) -- (top);
}
\end{tikzpicture}
\end{center}

The number of maximal chains in $\Pi_n$ can be counted by starting from $\hat{0} = \{1 \mid 2 \mid \ldots \mid n\}$ and then, at each level, joining two blocks from the partition.  
So the number of maximal chains in $\Pi_n$ is 
\begin{align*}
  \binom{n}{2} \binom{n - 1}{2} \cdots \binom{3}{2} \binom{2}{3}  
  & = \left( \frac{n!}{2! \cancel{(n - 2)!}} \right) \left( \frac{(n - 1)!}{2! \cancel{(n - 3)!}} \right) \cdots \left( \frac{\cancel{3!}}{2!} \right) \left( \frac{\cancel{2!}}{2!} \right) \\ 
  & = \frac{n! (n - 1)!}{2^{n - 1}}
\end{align*}

\section{}
Take any $\pi' \in [\hat{0}, \pi]_{\Pi_n}$. 
Since $\pi'$ refines $\pi$, there is a partition of $\pi'$ such that the union of each block in that partition is equal to a single block of $\pi$. 
In other words, if $\pi' = \{B_{1}' \mid B_{2}' \mid \ldots \mid B_{t}'\}$, then there is a reordering of $B_1, \ldots, B_k$ and $B'_{1}, \ldots, B'_{t}$ and numbers $1 < m_1 <  \ldots < m_{k - 1} < t$ such that 
$$B_1 = \bigsqcup_{i = 1}^{m_1} B_{i}', \: B_2 = \bigsqcup_{i = m_{1} + 1}^{m_2} B_{i}', \ldots, \:  B_k = \bigsqcup_{i = m_{k - 1} + 1}^{t} B_{i}' $$
Let $Q = \Pi_{\#B_1} \times \ldots \Pi_{\#B_{k}}$ and define $\varphi: [\hat{0}, \pi]_{\Pi_n} \rightarrow Q$ where 
\begin{equation} \label{eq:ptn_Bp}
  \pi' = \{B'_{1} \mid \ldots \mid B'_{t}\} \mapsto (\{B'_{1} \mid \ldots \mid B'_{m_1}\}, \ldots, \{B'_{m_{k - 1} + 1} \mid \ldots \mid B'_{t}\})
\end{equation}
Of course, the elements of the blocks of the partition in each coordinate of the image, for example $\{B'_{1} \mid \ldots \mid B'_{m_1} \}$, need to be relabeled to elements of $[\#B_1]$ in an order-preserving way. 
Then it is not difficult to see from the definition of $\varphi$ that given $\pi_1, \pi_2 \in [\hat{0}, \pi]_{\Pi_n}$, 
$$\varphi(\pi_1) \leq_{Q} \varphi(\pi_2) \iff \pi_1 \leq_{\Pi_n} \pi_2$$
so $\varphi$ is a homomorphism between posets. 
Given a refinement $\pi' \leq_{\Pi_n} \pi$, it is also clear that the partition of $\pi'$ given by \eqref{eq:ptn_Bp} is indeed unique to $\pi'$; 
no other $\pi'' \in [\hat{0}, \pi]_{\Pi_n}$ induces such a partition of $\pi'$. 
Finally, it is similarly trivial to see that $\varphi$ is surjective since a refinement of $\pi$ is essentially a refinement of each of the blocks of $\pi$. 

\section{}


\section{}
Notice $\text{Int}(P + Q) = \left\{ [x, y]_{P + Q} \mid (x, y \in P \text{ and } x \leq_{P}) \text{ or } (x, y \in Q \text{ and } x \leq_{Q} y) \right\}$. 
Hence 
$$\text{Int}(P + Q) \equiv \text{Int}(P) \sqcup \text{Int}(Q)$$
and the following functions can be defined: 
\begin{align*}
  \iota_{P}: \text{Int}(P) \hookrightarrow \text{Int}(P + Q) & ; \: [x, y]_{P} \mapsto [x, y]_{P + Q} \\ 
  \iota_{Q}: \text{Int}(Q) \hookrightarrow \text{Int}(P + Q) & ; \: [x, y]_{Q} \mapsto [x, y]_{P + Q} \\ 
\end{align*}

Take $f \in I(P + Q)$. 
Then define 
\begin{align*}
  g: \text{Int}(P) \rightarrow \C & ; \: [x, y]_{P} \overset{\iota_{P}}{\mapsto} [x, y]_{P + Q} \overset{f}{\mapsto} \C \\ 
  h: \text{Int}(Q) \rightarrow \C & ; \: [x, y]_{Q} \overset{\iota_{Q}}{\mapsto} [x, y]_{P + Q} \overset{f}{\mapsto} \C
\end{align*} 
Then if 
\begin{equation} \label{eq:oplus_map}
  (g, h)[x, y]_{P + Q} \defeq 
  \begin{cases}
    g[x, y]_{P} & \text{ if } x, y \in P \\ 
    h[x, y]_{Q} & \text{ if } x, y \in Q
  \end{cases}
\end{equation}
then it's clear $f = (g, h)$ and thus $f \in I(P) \oplus I(Q)$; 
this shows $I(P + Q) \subset I(P) \oplus I(Q)$. 

The reverse inclusion is clear: 
if $(g, h) \in I(P) \oplus I(Q)$, then $(g, h)$ can be viewed as an element in $I(P + Q)$ by defining $(g, h): \text{Int}(P + Q) \rightarrow \C$ as in \eqref{eq:oplus_map}. 

\section{}
Recall that $\zeta^2[x, y]$ counts the number of multichains $\{x \leq z \leq y\}$. 
Notice that these multichains, namely, via the middle element $z$, corresponds bijectively to an element in $[x, y]$. 
Hence $\zeta^2[x, y] = \#[x, y]$. 

\section{}
If $\hat{0} \lessdot x$, then 
$$0 = \sum_{\hat{0} \leq z \leq x} \mu[\hat{0}, z] = \mu[\hat{0}, \hat{0}] + \mu[\hat{0}, x] = 1 + \mu[\hat{0}, x]$$
hence $\mu[\hat{0}, x] = -1$. 
Next, suppose $y \lessdot x$ and $y \neq 0$. 
Then if $z \in P$ is an element such that $\hat{0} \leq z \leq x$ and $z \neq x$, then $z \leq y$. 
Therefore, 
\begin{align*}
  0 & = \sum_{\hat{0} \leq z \leq x} \mu[\hat{0}, z] \\ 
    & = \sum_{\hat{0} \leq z \leq y} \mu[\hat{0}, z] + \mu[\hat{0}, x] \\ 
    & = \mu[\hat{0}, x]
\end{align*}
as was to be shown. 

\section{}
For $n = 0$, let $P$ be a chain $\{x_1 < x_2 < x_3\}$. 
Then $x_1 = \hat{0}$ and 
$$0 = \sum_{\hat{0} \leq z \leq x_3} \mu_[\hat{0}, z] \mu[\hat{0}, x_3] = 1 - 1 + \mu[\hat{0}, x_3] = \mu[\hat{0}, x_3]$$

Next, take $n > 0$. 
Then let $P$ be a graded poset of rank 2, with minimum $\hat{0}$ and maximum $\hat{1}$,
i.e., the only element of rank 0 is $\hat{0}$ and the only element of rank 2 is $\hat{1}$. 
Finally, let there be $n + 1$ elements of rank 1. 
For all $z \in P$ with $\rank{z} = -1$, then $\mu[\hat{0}, z] = 1$. 
Then 
\begin{align*}
  0 & = \sum_{\hat{0} \leq z \leq \hat{1}} \mu[\hat{0}, z] \\ 
    & = \mu[\hat{0}, \hat{0}] + \sum_{\substack{z \in P \\ \rank{z} = 1}} \mu[\hat{0}, z] + \mu[\hat{0}, \hat{1}] \\ 
    & = 1 + -(n + 1) + \mu[\hat{0}, \hat{1}] 
\end{align*}
Hence $\mu[\hat{0}, \hat{1}] = n$. 

Finally, let $n < 0$. 
If $n$ is even, then let $P$ be a poset of rank 3 with minimum $\hat{0}$ and $\hat{1}$. 
Add three elements of rank 1 and $\abs{n} / 2 + 1$ elements of rank 2, where each element of rank 2 is connected to all three elements of rank 1. 
Of course, for every $z \in P \with \rank{z} = 1$, $\mu[\hat{0}, z] = -1$. 
Furthermore, for every $z \in P$ with $\rank z = 2$, $\mu[\hat{0}, z] = 2$. 
Hence 
\begin{align*}
  0 & = \sum_{\hat{0} \leq z \leq \hat{1}} \mu[\hat{0}, z] \\ 
    & = \mu[\hat{0}, \hat{0}] + \sum_{\substack{z \in P \\ \rank{z} = 1}}\mu[\hat{0}, z] + \sum_{\substack{z \in P \\ \rank{z} = 2}}\mu[\hat{0}, z] + \mu[\hat{0}, \hat{1}] \\ 
    & = 1 + 3(-1) + (\frac{\abs{n}}{2} + 1)2 + \mu[\hat{0}, \hat{1}] \\ 
    & = \frac{\abs{n}}{2} + \mu[\hat{0}, \hat{1}] 
\end{align*}
Hence $\mu[\hat{0}, \hat{1}] = -\abs{n} = n$ since $n < 0$.  

\newpage{}
\section*{Acknowledgement}
I used claude.ai to typeset the Hasse Diagrams of $\Pi_3$ and $\Pi_4$ after first illustrating them for myself on paper. 

\end{document}
